\documentclass[11pt]{article}
\usepackage{amsmath,amssymb,booktabs}
\usepackage[margin=1in]{geometry}

\title{MVP Simulator: MDP Formulation and Initial Results}
\author{William Dennis}
\date{June 2026}

\begin{document}
\maketitle

\section{MDP Formulation}

The MVP simulator implements a finite-horizon MDP $(\mathcal{S}, \mathcal{A}, P, r, \gamma, T)$:

\subsection{State Space}
\[
\mathcal{S} = \{\texttt{sedentary},\ \texttt{active}\}
\]
Two activity states with no additional time-of-day dimension. Time-dependent reward scaling is handled via \texttt{reward\_multiplier\_by\_step} (see \S\ref{sec:reward} below).

\subsection{Action Space}
\[
\mathcal{A} = \{\texttt{nudge},\ \texttt{idle}\}
\]
A binary intervention: send a motivational message (\texttt{nudge}) or do not (\texttt{idle}).

\subsection{Transition Function}

$P(s' \mid s, a)$ is defined by a configurable matrix. The probabilities encode a Reactance effect: reminders help sedentary users but can backfire for users who are already active.

\begin{table}[h]
\centering
\begin{tabular}{llcc}
\toprule
Current state $s$ & Action $a$ & $P(\texttt{active})$ & $P(\texttt{sedentary})$ \\
\midrule
sedentary & nudge & 0.3 & 0.7 \\
sedentary & idle & 0.1 & 0.9 \\
active & nudge & 0.5 & 0.5 \\
active & idle & 0.6 & 0.4 \\
\bottomrule
\end{tabular}
\caption{Transition probabilities. Sending a reminder (nudge) increases the probability of becoming active for sedentary users ($0.3 > 0.1$) but decreases it for active users ($0.5 < 0.6$), consistent with Reactance theory.}
\end{table}

Optional \texttt{reward\_multiplier\_by\_step}: a per-step scaling factor that can zero out rewards at certain daily steps (soft mask). When absent, all multipliers default to 1.0. The per-step reward is computed as $\text{base reward} \times \text{multiplier}[t]$ where $t$ is the step-of-day index.

\subsection{Reward}\label{sec:reward}
\[
r(s, a, t) = \texttt{per\_step\_reward}[t][s']
\]
where $s'$ is the next state after taking action $a$ in state $s$, and $t \in \{0, \dots, \texttt{steps\_per\_day} - 1\}$
is the step-of-day index. The precomputed table is:
\[
\texttt{per\_step\_reward}[t][s] = \texttt{reward}[s] \times \texttt{reward\_multiplier\_by\_step}[t]
\]
with a default multiplier of 1.0 for all $t$ when not specified. Currently:
\begin{align*}
\texttt{reward}[\texttt{active}] &= 1.0 \\
\texttt{reward}[\texttt{sedentary}] &= 0.0
\end{align*}
Single-timescale placeholder. Multi-timescale reward (immediate + delayed body measure) is deferred to Phase 2.

\subsection{Episode Structure}
$T = \texttt{episode\_days} \times \texttt{steps\_per\_day}$ timesteps. Default: $T = 90 \text{ days} \times 5 \text{ steps/day} = 450$.

\section{Agents}

All agents are contextual bandits: state is accepted but not used in action selection. For the stationary MVP transition matrix, $E[r \mid a]$ is the relevant quantity.

\subsection{Thompson Sampling}
Beta-Bernoulli TS with prior $\text{Beta}(1, 1)$ per action. Posterior update:
\[
\alpha_a \leftarrow \alpha_a + \mathbb{1}[s'.\texttt{activity} = \texttt{active}], \quad
\beta_a \leftarrow \beta_a + \mathbb{1}[s'.\texttt{activity} = \texttt{sedentary}]
\]
Action selected by sampling $\theta_a \sim \text{Beta}(\alpha_a, \beta_a)$ and choosing $a^* = \arg\max_a \theta_a$.

\subsection{Epsilon-Greedy}
$\epsilon = 0.1$. With probability $\epsilon$, select uniformly at random; otherwise select $\arg\max_a Q(a)$. Q-values updated via incremental average:
\[
Q(a) \leftarrow Q(a) + \frac{r_t - Q(a)}{n_a}
\]

\subsection{UCB1}
Upper Confidence Bound with exploration parameter $c = 2.0$. Selects:
\[
a^* = \arg\max_a \left[ Q(a) + c \sqrt{\frac{\ln N}{N_a}} \right]
\]
where $N$ is total steps and $N_a$ is the number of times action $a$ was chosen. The confidence bonus shrinks as $N_a$ grows, naturally shifting from exploration to exploitation. Each action is pulled once before UCB is applied.

\subsection{Random Baseline}
Uniform random action selection. No learning. Serves as lower bound.

\section{Initial Results}

All agents run over 10 random seeds (1--10), 450 timesteps each. Results are reported as mean $\pm$ standard deviation across seeds.

\subsection{Standard Config (\texttt{config/rule\_based.yaml})}
No \texttt{reward\_multiplier\_by\_step} -- uniform 1.0 multiplier on all steps.

\begin{table}[h]
\centering
\begin{tabular}{lrr}
\toprule
Agent & Total Reward & Mean Reward/Step \\
\midrule
Thompson Sampling ($\alpha_0=\beta_0=1$) & 160.3 $\pm$ 10.5 & 0.356 $\pm$ 0.023 \\
Epsilon-Greedy ($\epsilon=0.1$) & 160.8 $\pm$ 15.3 & 0.357 $\pm$ 0.034 \\
UCB ($c=2.0$) & 148.2 $\pm$ 15.0 & 0.329 $\pm$ 0.033 \\
Random & 135.6 $\pm$ 10.7 & 0.301 $\pm$ 0.024 \\
\bottomrule
\end{tabular}
\caption{Results on \texttt{config/rule\_based.yaml} (10 seeds, 450 timesteps). TS and $\epsilon$-greedy perform similarly; TS has lower variance. UCB outperforms random but shows higher variance.}
\end{table}

\subsection{Masked Config (\texttt{config/rule\_based\_with\_mask.yaml})}
\texttt{reward\_multiplier\_by\_step: [1, 1, 1, 1, 0]} -- step 4 (end of day) yields zero reward regardless of outcome. All agents drop proportionally to the masked fraction ($1/5$ of steps zeroed), confirming the multiplicative reward model.

\begin{table}[h]
\centering
\begin{tabular}{lrr}
\toprule
Agent & Total Reward & Mean Reward/Step \\
\midrule
Thompson Sampling ($\alpha_0=\beta_0=1$) & 126.6 $\pm$ 6.8 & 0.281 $\pm$ 0.015 \\
Epsilon-Greedy ($\epsilon=0.1$) & 127.3 $\pm$ 10.2 & 0.283 $\pm$ 0.023 \\
UCB ($c=2.0$) & 116.9 $\pm$ 8.6 & 0.260 $\pm$ 0.019 \\
Random & 107.8 $\pm$ 6.7 & 0.240 $\pm$ 0.015 \\
\bottomrule
\end{tabular}
\caption{Results on \texttt{config/rule\_based\_with\_mask.yaml} (10 seeds, 450 timesteps). Reward drops $\approx 20\%$ across all agents, matching the $1/5$ masked fraction. Relative ordering is preserved.}
\end{table}

\section{Known Limitations}

\begin{itemize}
\item Agents are state-blind bandits, not state-aware RL. State-aware agents are Phase 2.
\item Single-timescale reward only. Multi-timescale (immediate + 3-week delayed) is Phase 2.
\item No user response model. Burden accumulation and 4 archetypes are Phase 2.
\item Transition matrix is fixed and known. Learned transitions from real data are Phase 2.
\item Transition probabilities are placeholder values informed by Reactance theory. Calibration against real data is Phase 2.
\end{itemize}

\end{document}
